\section{Discussion}\label{sec:discussion}

Several parts of the system are not mechanised. The mechanised tally decrypts with a single trustee key. In a
deployment the key is shared: trustee $i$ holds $x_i$ and publishes $h_i =
g^{x_i}$ with a Schnorr proof of knowledge, which prevents a rogue-key
attack, and the election key is $h = h_1 \cdots h_n$. Ballots, their
proofs, and aggregation use only $h$, so their theorems hold unchanged.
Only part (i) of the certificate uses the secret key. With shared keys it
becomes one partial decryption per trustee, each with a Chaum--Pedersen
proof against $h_i$. The library's \rocq{DistElgamal} module proves that
the partial decryptions combine correctly. Parts (ii)--(iv) need only the
randomness of the fresh encryptions and do not change. The lemma that $n$
accepted partial-decryption proofs imply a correct combined decryption is
not yet mechanised. The sharing is $n$-of-$n$, so privacy
needs one honest trustee but every trustee must take part. A $t$-of-$n$
sharing would add Shamir shares and Lagrange interpolation.
The Fiat--Shamir step is on paper: the
non-interactive theorems (\rocq{nizk_ballot_valid},
\rocq{nizk_tally_complete}) are completeness statements, and
soundness is stated in the rewinding form of special soundness.
Connecting the library to a formalisation of the random oracle
model~\cite{10.1007/11617990_3} is future work. Ballot privacy and
receipt-freeness combine the zero-knowledge theorems with an on-paper
IND-CPA reduction (Section~\ref{sec:security}), and there is no game-based
receipt-freeness theorem. The closure algorithm of the honest trustees
is unverified, but soundness does not depend on it, because its certificate is
checked.

The theorems of Section~\ref{sec:mech} are about what the certified
verifier checks. Unverified code decides what reaches it and about which
values, and the Belenios defects were in such code. The security of the
scheme depends on three decisions in that code. The decoder that admits a
ballot from the board must check that every group element has order $q$
and every field element is reduced, since \rocq{ballot_sound} is a
theorem about elements of the group type. The verifier must \emph{derive}
the $\delta$ ciphertexts from the position ciphertexts and never accept
them as inputs (Section~\ref{sec:ballot}). A driver that read
$\mathrm{Enc}(\delta_{ij})$ from the ballot would verify every proof and
count a ballot twice. Finally, the Fiat--Shamir challenge of a
ballot must bind the voter's credential, or a ballot can be replayed
under another one. The mechanised theorems hold for any oracle. The
driver instantiates the oracle per voter (Section~\ref{sec:eval:gentoo}), and
a driver that did not would be wrong. Every run of the Gentoo driver
tests the credential binding: it replays a ballot under another
credential and checks that the verifier rejects it. The types enforce
the derivation of $\delta$: the mechanised ballot has no field for
$\mathrm{Enc}(\delta_{ij})$, so no driver can supply one, and
\rocq{verify_ballot} computes it. This wiring is protocol logic that no
current formalisation covers end to end. The \rocq{count} judgement
covers the aggregation and the final check. The next proof we would
write extends it to the derivation of the statements each proof is
about.

The number of certificate proofs is independent of $N$, but each range
proof and each closure disjunction is a membership proof with $O(N)$
branches, so the certificate has $O(m^3 N)$ group elements and the
trustees' and verifier's time grows linearly with the electorate
(Section~\ref{sec:eval:tally}). The first ballot has $O(m^3)$ group
elements for the same reason, $O(m)$ branches in each of $m^2$ cells.
The batched ballot of Section~\ref{sec:batched} removes the branches
and costs $O(m^2 \log m)$. The same batching applies to the certificate.
The trustees know the openings of $d$ and $p$, so every range check of
Section~\ref{sec:tally} is a bit-decomposed range of a linear form of
committed vectors, provable by one zero argument, and the closure
disjunction becomes a committed minimum with one product check. The
certificate would then be $O(m^3 \log N)$ field elements instead of
$O(m^3 N)$ group elements. We have not implemented this.

The witness paths are the certificate's only leakage beyond the winner
(Section~\ref{sec:security}), and a certificate can do without them.
Starting from $\mathrm{Enc}(d)$, $\lceil \log_2 (m-1) \rceil$ rounds of
encrypted min-max squaring $p' = \max(p, p \otimes p)$ reach the closure.
Each entry of each round needs $p' \ge p$ and the closure disjunctions of
(iii) for $p' \ge \min(p_{ik}, p_{kj})$, and, for $p'$ to be no larger,
one $(m{+}1)$-way disjunction stating that $p'$ is at most $p_{ij}$ or
at most both $p_{ik}$ and $p_{kj}$ for some $k$. This gives $O(m^3 \log m)$
disjunctions and no cleartext beyond the winner. We have neither
implemented nor mechanised it.

The two ballots rest on different assumptions
(Table~\ref{tab:twoballots}). The first is sound by exact special
soundness from two transcripts, with no assumption beyond the election
group. The batched ballot assumes binding of the Pedersen commitment key,
whose generators are hashed to the group as Verificatum and Swiss Post
do. Its soundness is statistical and multi-round, in the
perfect-adversary form of the library's shuffle theorems, with the
forking argument and its Fiat--Shamir counterpart~\cite{afk21} on paper.
Its zero-knowledge names the discrete logarithms of the key
(Section~\ref{sec:batched}). None of these assumptions is new to the
library or to deployed systems, but the first ballot makes none of them,
which is why the paper keeps both ballots.
